% 06-verifiers.tex: independent verification.
\section{Verification without the issuer: the detached verifier, the SDKs and the conformance contract}
\label{sec:verifiers}

\subsection{The detached verifier}

\code{scripts/polaris-verify.py} verifies a credential's authenticity pack with only a standard
ML-DSA library: no Polaris code, no database, no network. It is the one Polaris component whose
input is chosen by an adversary, because a relying party runs it on whatever a stranger presented,
and it is the component through which every later protocol artifact is checked: manifests, epoch
checkpoints, revocation feeds, status bundles, status assertions, transparency heads, receipts,
timestamps, signed documents, registries, trust lists, presentations. Three independent ML-DSA
implementations agree on the published vectors: a lattice library, OpenSSL, and a third in the
TypeScript SDK. \sImpl

The verifier is held \emph{total} against hostile input by a metamorphic fuzzer: for every signed
type it builds a genuine object, confirms acceptance, then flips bits in the signature and key,
mutates and drops every signed field, substitutes adversarial values, and feeds each object to every
other type's verifier, asserting a clean fail-closed rejection with no exception. The fuzzer found
that the verifier crashed on several classes of malformed input, an integer where a list was
expected, a hex field of the wrong type, and the fix was a totality contract that converts every
would-be crash into the reject the verifier already intended. A verifier that raises pushes the
fail-closed decision onto a caller who may not catch it. \sImpl

An attack suite runs beside it: executable adversaries that try to forge a signature with an
attacker key, tamper a byte, present a signature against a different token, relabel the placeholder
as real, or present a revoked token, against the real code. An attack succeeds when a defence fails,
and the build goes red if any succeeds. A grep proves a line exists; an attack proves a defence holds.
\sImpl

\subsection{The SDKs and the conformance contract}

A relying party installs one of two verify SDKs, a Python reference and a TypeScript implementation,
each answering one narrow question about a presented credential, is it authentic and is it
authoritative right now, and never returning a person's data. Both are held to a language-agnostic
conformance suite: a runner drives any verifier over standard input and output through the published
cases, forty-eight at \code{v9.345} across sixteen artifact types, including the composite federation
trust decision. Passing the suite is what \emph{conformant} means, and the normative wire
specification (RFC 2119 language, twenty format strings, the exact signed-field list of every
artifact) is what an implementation importing no Polaris code would build against. \sImpl

That last sentence is the honest boundary of this layer. The specification, the cases and two SDKs
exist, and a check pins the specification to the signer so a documented field list cannot drift from
the code. An actual implementation by an external team, from the documents alone, has not happened;
the roadmap marks it as blocked on an external actor. The conformance suite is the contract. The
integration is unproven. \sExt

\subsection{Versioning, and the frozen version 1}

Protocol majors live in the format string (\code{polaris-registry/1}), minors are advertised by the
registry and are never needed to verify, and the negotiation rule is normative: reject an unknown
major, accept any minor, never emit an unadvertised version, answer an unsupported one with a
specific error and the supported list. Version 1 of the protocol is frozen under checksums, with the
detached verifier of \code{v9.317} vendored beside it, and a compatibility suite runs both directions
on every push: today's verifiers must accept everything version 1 published, and the pinned older
verifier must never accept what today's suite rejects. \sImpl

\begin{layercard}{Layer card: independent verification}
\qa{What it is}{A standalone verifier, two SDKs, a conformance runner over published cases, a normative wire specification, a frozen version-1 set with a cross-version suite, a fuzzer and an attack suite.}
\qa{Why it exists}{A verification the issuer performs asks the relying party to trust the issuer's server. Authenticity has to be decidable by anyone, anywhere, with the network off.}
\qa{What it trusts}{The published issuer keys the relying party chose to trust (its anchor set); a standard ML-DSA library.}
\qa{What it refuses}{Any Polaris code or database; any input it cannot parse; any object presented to the wrong type's verifier; any version it was not told about.}
\qa{What it can see}{The artifact and the anchor set. Nothing else.}
\qa{What it cannot see}{Whether the credential is still authorized, unless a status artifact is presented with it.}
\qa{If it fails}{A wrongly accepted mutation fails the fuzzer; a defence that no longer holds fails the attack suite; a drift between specification and signer fails \code{check\_wire\_spec\_matches\_code}; a frozen case that stops verifying fails the compatibility suite.}
\qa{Checked by}{\code{check\_detached\_verifier}, \code{check\_conformance\_suite}, \code{check\_typescript\_sdk}, \code{check\_verifier\_fuzz}, \code{check\_attacks\_run}, \code{check\_frozen\_v1}.}
\qa{Now or later}{Now, for the software. The external implementation is later and not ours to do.}
\qa{Inside and outside}{Inside: the signatures. Outside: the status layer, which answers the question a signature cannot.}
\end{layercard}

\nextlayer{Signing a credential proves who signed it. It does not prove whether that credential is
still authorized: the seal on a revoked passport is as genuine as the seal on a valid one. Polaris
therefore needs status.}
